\begin{multicols}{2}
\section{Prompt for Test Refactoring and GroundPRD Construction}
\label{appendix:task1-prompt}
\end{multicols}

\begin{promptbox}[GroundPRD Construction \& Test Refactoring]
\textbf{\#\# Task Description}

You are given an original runnable Docker environment containing the repository, its runtime environment, and its original test scripts. Without modifying the repository's functional code, refactor the original tests into a standardized black-box input/output format and produce a PRD that describes only ``what input leads to what output.''

\textbf{\#\# Inputs}

Original Docker install and test script: \texttt{\{test\_script\}}

\textbf{\#\# Outputs}
\begin{itemize}
\item \texttt{repos/\{username\}\_\_\{repo\_name\}/rcb\_tests/test\_cases/featureN\_[name].json} --- a subset of original repository tests used as the hidden evaluation set
\item \texttt{repos/\{username\}\_\_\{repo\_name\}/rcb\_tests/public\_test\_cases/featureN\_[name].json} --- the strict mirror of the example cases embedded in the PRD
\item \texttt{repos/\{username\}\_\_\{repo\_name\}/start.md}
\item \texttt{repos/\{username\}\_\_\{repo\_name\}/rcb\_tests/test.sh} --- a single entry point supporting \texttt{--cases-dir <subdir>}
\item \texttt{repos/\{username\}\_\_\{repo\_name\}/<dispatcher source>} --- placed idiomatically for the chosen language
\end{itemize}

\textbf{\#\# Workflow}
\begin{enumerate}
\item Read the source code and the original repository tests.
\item For each original test, decide whether it checks externally visible behavior that is independent of the implementation choice. Keep those tests; discard white-box or implementation-specific assertions.
\item Write the PRD and generate \texttt{test\_cases}: each retained test is converted into JSON cases, and the PRD describes the behavior in pure input/output terms. The PRD must be self-contained: each leaf feature block should be understandable from its text and embedded example cases alone.
\item Run \texttt{bash rcb\_tests/test.sh --cases-dir test\_cases} inside the original Docker image. The result must be zero FAIL. If a case fails, either the dispatcher is inconsistent with the golden implementation or the case should not have been retained.
\end{enumerate}

\textbf{\#\# Requirements}
\begin{enumerate}
\item Do not modify the repository's functional code. The only files you may add or change are the test data, \texttt{test.sh}, and the Task-1 dispatcher/adapter.
\item \texttt{bash rcb\_tests/test.sh} must support \texttt{--cases-dir <subdir>} (default \texttt{test\_cases}). For every case it writes the raw stdout to \texttt{rcb\_tests/stdout/<cases-dir>/\{stem\}@\{idx:03d\}.txt}. Different case directories must never overwrite one another.
\item Each JSON file uses schema \texttt{\{description, cases:[\{input, expected\_output\}]\}}. One leaf feature corresponds to one JSON file.
\item \texttt{test\_cases} must remain a subset of the original repository tests. Do not invent new behavior at this stage.
\item \texttt{input} must be the program's real input data, not source-level method names or class names.
\item \texttt{expected\_output} must be the program's real stdout and must be directly comparable to the contents of \texttt{rcb\_tests/stdout/*.txt}. Avoid weak pass/fail markers or single booleans when richer behavior can be exposed.
\item Errors must be normalized into language-neutral black-box contracts rather than leaking host-language exception names or runtime-specific formatting.
\item \texttt{description} must let a reader understand the scenario together with \texttt{input} and \texttt{expected\_output}, but it must not copy the concrete input values or exact expected output verbatim.
\item The PRD must not leak original repository identifiers or implementation details such as repository names, author names, original class or method names, namespaces, or language-locked framework labels.
\item Use multi-level features when necessary. Every leaf feature corresponds to one independent functionality point and one JSON file.
\end{enumerate}

\textbf{\#\# PRD Template}

The PRD should contain a neutral function-based title, a project goal, background and problem statement, architecture and engineering constraints, and a core-features section in which each leaf feature includes: a short user-story style statement, expected behavior in plain language, a pointer to \texttt{rcb\_tests/public\_test\_cases/...}, and one or more embedded example JSON cases.

\textbf{\#\# Deliverables}
\begin{enumerate}
\item A maintainable core system description reflected in the PRD.
\item An execution/test adapter that reads JSON from stdin, invokes the core behavior, and prints the required stdout contract.
\item An automated harness \texttt{bash rcb\_tests/test.sh} that runs all cases from a selected case directory and records per-case stdout files exactly as specified.
\end{enumerate}
\end{promptbox}
